\section{Problem Statement and Hardware-Software System Architecture}

\subsection{Operational Challenges}
The deployment of IoT-Based Smart Health Monitoring System in operational environments faces Motion artifacts corrupting optical PPG waveforms, strict battery life constraints on wearable wristbands, biometric data privacy under HIPAA/GDPR, and latency in dispatching emergency alerts.. The complete system architecture must maintain high packet delivery reliability, protect data in transit via cryptographic primitives, and deliver deterministic response latency.

\subsection{Hardware and Cloud Pipeline Architecture}
Figure~\ref{fig:pipeline} presents the multi-tier architectural pipeline implemented for this case study. It partitions the data lifecycle across Perception, Edge Gateway, Network Transport, and Cloud Analytics layers.

\begin{figure}[H]
\centering
\resizebox{\textwidth}{!}{\begin{tikzpicture}[
    block/.style={rectangle, draw=black!80, fill=green!8, thick, rounded corners, minimum height=1.1cm, text width=2.5cm, font=\footnotesize},
    arr/.style={-{Stealth[length=2mm]}, thick}
]
\node[block, align=center] (n1) at (0,0) {1. Patient Wearable\\MAX30102 \& MLX90614};
\node[block, align=center] (n2) at (3.2,0) {2. Biosignal Filtering\\Bandpass \& Motion Rejection};
\node[block, align=center] (n3) at (6.4,0) {3. Micro-Controller\\Nordic nRF52 BLE 5.0 SoC};
\node[block, align=center] (n4) at (9.6,0) {4. AES-128 Encryption\\Crypto-Hardware Engine};
\node[block, align=center] (n5) at (12.8,0) {5. Mobile Gateway\\Patient Smartphone Relayer};
\node[block, align=center] (n6) at (12.8,-2.2) {6. Hospital Cloud\\FHIR-Compliant Server};
\node[block, align=center] (n7) at (9.6,-2.2) {7. Real-Time Telemetry\\Kafka Stream Analytics};
\node[block, align=center] (n8) at (6.4,-2.2) {8. Anomaly Detector\\Tachycardia/Arrhythmia ML};
\node[block, align=center] (n9) at (3.2,-2.2) {9. Emergency Dispatch\\Automated SMS/Push to Doc};
\node[block, align=center] (n10) at (0,-2.2) {10. Clinical Dashboard\\Doctor Portal \& Triage};
\draw[arr] (n1) -- (n2);
\draw[arr] (n2) -- (n3);
\draw[arr] (n3) -- (n4);
\draw[arr] (n4) -- (n5);
\draw[arr] (n5) -- (n6);
\draw[arr] (n6) -- (n7);
\draw[arr] (n7) -- (n8);
\draw[arr] (n8) -- (n9);
\draw[arr] (n9) -- (n10);
\end{tikzpicture}}
\caption{End-to-End System Architecture and Telemetry Pipeline for IoT-Based Smart Health Monitoring System.}
\label{fig:pipeline}
\end{figure}

\subsection{Architectural Component Specifications}
\begin{enumerate}
  \item \textbf{Perception Layer:} Industrial-grade transducers interface directly with the physical domain, utilizing differential analog front-ends or SPI/I2C digital buses.
  \item \textbf{Edge Microcontroller Node:} Low-power 32-bit Arm Cortex or ESP32 processor executes local FIR filtering, sensor calibration lookup tables, and payload serialization.
  \item \textbf{Secure Gateway Layer:} Aggregates wireless packets from distributed end-nodes, performs protocol translation (e.g., LoRaWAN to MQTT/TLS), and routes over cellular or fiber uplinks.
  \item \textbf{Cloud Broker \& Storage:} Scalable clustered MQTT brokers stream ingress telemetry into time-series databases (InfluxDB) optimized for downsampling and high-write concurrency.
  \item \textbf{Decision Support UI:} Interactive Grafana dashboards, automated alerting microservices, and SCADA relays actuate physical interventions based on analytical threshold triggers.
\end{enumerate}
